\documentclass[11pt]{article}

\usepackage[T1]{fontenc}
\usepackage{amsmath,amssymb,amsthm,mathtools}
\usepackage{geometry}
\usepackage{graphicx}
\usepackage{microtype}
\usepackage{enumitem}
\usepackage[hidelinks]{hyperref}
\geometry{margin=1in}

\newtheorem{theorem}{Theorem}
\newtheorem{lemma}{Lemma}
\newtheorem{proposition}{Proposition}
\newtheorem{corollary}{Corollary}
\theoremstyle{remark}
\newtheorem{remark}{Remark}

\newcommand{\T}{\mathbb T}
\newcommand{\R}{\mathbb R}
\newcommand{\C}{\mathbb C}
\newcommand{\Z}{\mathbb Z}
\newcommand{\eps}{\varepsilon}
\newcommand{\cO}{\mathcal O}
\newcommand{\phg}{\mathrm{phg}}

\title{The Boutet de Monvel conjecture for strongly convex\\
Fourier multipliers on flat tori}
\author{Candidate partial result for arXiv:2110.09389\\
Run \texttt{fa\_banach\_001}, lane 10, model GPT5.6}
\date{21 July 2026}

\begin{document}
\maketitle

\begin{center}
\fbox{\parbox{0.91\linewidth}{
\textbf{Status.} Candidate partial result, likely valid, subject to expert review.
The theorem below verifies all six conclusions of Conjecture 4.1 for a class of
translation-invariant analytic elliptic operators on flat tori.  The class
contains fourth-order differential operators that are not powers of a
Laplacian.  The conjecture on general analytic manifolds remains open.
}}
\end{center}

\section{Source problem}

The source is D. S. Winterrose, \emph{Algebras of pseudo-differential
operators acting on holomorphic Sobolev spaces}, arXiv:2110.09389,
Conjecture 4.1 on printed page 8 \cite{Winterrose}.  For an analytic, formally
self-adjoint, nonnegative elliptic polyhomogeneous pseudodifferential operator
$P$ of order $d>0$, with positive principal symbol $p$ and strictly convex
unit co-balls, the conjecture asserts six conclusions.  In abbreviated form:

\begin{enumerate}[label=(\arabic*)]
\item the imaginary-time Hamilton flow of $p^{1/d}$ identifies
      $B^*_{\eps}M\setminus0$ with a complex tube
      $M^\Phi_\eps\setminus M$;
\item the tube is open with real-analytic boundary;
\item it is strictly pseudoconvex;
\item the kernel of $e^{-\eps P^{1/d}}$ extends holomorphically in its first
      variable to the tube;
\item its boundary value is a polyhomogeneous complex-phase Fourier integral
      operator of order $-(n-1)/4$;
\item that operator is a homeomorphism
      $H^s(M)\to\cO^{s+(n-1)/4}(\partial M^\Phi_\eps)$ for every $s\in\R$.
\end{enumerate}

Here $\cO^r(\partial D)$ is the closure in $H^r(\partial D)$ of smooth
boundary values of functions holomorphic in $D$, as in the source.  The full
source statement is reproduced in Figure~\ref{fig:problem}.

\begin{figure}[ht]
  \centering
  \includegraphics[width=0.98\linewidth]{figures/open_problem_crop.png}
  \caption{Conjecture 4.1 in arXiv:2110.09389, printed page 8.}
  \label{fig:problem}
\end{figure}

\section{Statement of the partial result}

Write $\T^n=(\R/2\pi\Z)^n$ and
$\T^n_{\C}=(\C/2\pi\Z)^n$.  Let
$q:\R^n\to[0,\infty)$ have the following properties:

\begin{enumerate}[label=(Q\arabic*)]
\item $q(0)=0$, while $q$ is positive, even, positively one-homogeneous, and
      real analytic on $\R^n\setminus0$;
\item for $F(\xi)=\tfrac12q(\xi)^2$, the Hessian $D^2F(\xi)$ is positive
      definite for every $\xi\ne0$.
\end{enumerate}

Thus $q$ is a strongly convex analytic Minkowski norm.  Its dual norm is
\[
q^*(v)=\sup_{q(\xi)=1}v\cdot\xi .
\]
For $d>0$, define the Fourier multiplier $P=P_{q,d}$ by
\[
P e_k=q(k)^d e_k,
\qquad e_k(x)=e^{ik\cdot x},\quad k\in\Z^n,
\]
with $q(0)^d=0$.  Equivalently, modulo a smoothing modification at the zero
section, $P$ has analytic classical symbol $p(\xi)=q(\xi)^d$.

\begin{theorem}[All six conclusions on a flat torus]\label{thm:main}
Let $q$ satisfy \textnormal{(Q1)--(Q2)} and let $P=P_{q,d}$.  Then for every
$\eps>0$, all six conclusions of Winterrose's Conjecture 4.1 hold.  More
precisely,
\[
D_\eps:=\{v\in\R^n:q^*(v)<\eps\},
\qquad
M^\Phi_\eps=\T^n+iD_\eps,
\]
and the flow diffeomorphism is
\[
\Phi(x,\xi)=x+i\nabla F(\xi).
\]
The boundary operator is
\begin{equation}\label{eq:S}
(S_\eps f)(x+iv)
=\sum_{k\in\Z^n}e^{-\eps q(k)}\widehat f(k)
  e^{ik\cdot x-k\cdot v},
\qquad q^*(v)=\eps,
\end{equation}
and it is a complex-phase polyhomogeneous FIO of order $-(n-1)/4$ satisfying
\[
S_\eps:H^s(\T^n)\xrightarrow{\ \sim\ }
\cO^{s+(n-1)/4}(\partial M^\Phi_\eps)
\quad(s\in\R).
\]
In fact the admissible radii are unbounded: one may take
$\eps_0=\eps'_0=+\infty$.
\end{theorem}

\begin{corollary}[A non-Laplacian differential example]\label{cor:quartic}
There is $\delta_0>0$ such that, for every $0<\delta<\delta_0$, the operator on
$\T^2$
\[
P_\delta=(-\Delta)^2+\delta\,\partial_1^2\partial_2^2
\]
satisfies all six conclusions of Conjecture 4.1 for every $\eps>0$.  Its
principal symbol
\[
p_\delta(\xi)=(\xi_1^2+\xi_2^2)^2+\delta\xi_1^2\xi_2^2
\]
is not the square of a quadratic form when $\delta\ne0$; hence this is not the
flat-Laplacian case in disguise.
\end{corollary}

\section{Proof intuition}

On a translation-invariant torus the Hamiltonian covector never moves.  At the
imaginary time prescribed in the conjecture, its base point moves by
$i\nabla(\tfrac12q^2)$.  Legendre duality identifies this gradient map exactly
with the dual-norm tube $q^*(\operatorname{Im}z)<\eps$.  The same duality gives
the sharp estimate that makes the Fourier series of the Poisson kernel
holomorphic in the tube.

On the boundary, Poisson summation turns the Fourier series into the positive
complex phase
\[
(x-y)\cdot\xi+i\bigl(\eps q(\xi)+v\cdot\xi\bigr).
\]
The imaginary part is nonnegative by duality and vanishes precisely at the
unique support point $v=-\eps\nabla q(\xi)$.  Strong convexity makes this a
nondegenerate positive phase.  The standard order formula gives
$-(n-1)/4$.  Finally, Laplace's method on the dual unit sphere shows that
$S_\eps^*S_\eps$ is elliptic of order $-(n-1)/2$.  This gives the exact
Sobolev gain; Fourier polynomials and closed range then identify the range
with the whole holomorphic Sobolev space.

The mechanism relies decisively on translation invariance.  It does not prove
the conjecture for variable symbols or curved manifolds.

\section{Proof}

\subsection{Legendre geometry and the imaginary flow}

\begin{lemma}[Gradient parametrization of the dual ball]\label{lem:legendre}
Let $F=\tfrac12q^2$ and let $F^*$ be its Fenchel dual.  Then
\[
F^*(v)=\tfrac12q^*(v)^2.
\]
Moreover, $\nabla F:\R^n\setminus0\to\R^n\setminus0$ is a real-analytic
diffeomorphism with inverse $\nabla F^*$, and
\[
q^*(\nabla F(\xi))=q(\xi).
\]
Consequently $\nabla F$ maps $\{0<q<\eps\}$ real-analytically and
diffeomorphically onto $\{0<q^*<\eps\}$.
\end{lemma}

\begin{proof}
The positive Hessian and two-homogeneity of $F$ make $F$ strictly convex and
coercive.  Hence its gradient is injective and proper away from the origin,
and the inverse-function theorem makes it a global real-analytic
diffeomorphism of the punctured spaces.  Equality in Fenchel duality gives
$\nabla F^*=(\nabla F)^{-1}$.  Maximizing
$v\cdot\xi-\tfrac12q(\xi)^2$ first along every ray gives
$F^*(v)=\tfrac12q^*(v)^2$.

Euler's identity gives
$\xi\cdot\nabla F(\xi)=2F(\xi)=q(\xi)^2$.  Equality in Fenchel duality and
one-homogeneity then give
$q^*(\nabla F(\xi))=q(\xi)$.  The last assertion follows.
\end{proof}

The Hamiltonian of the conjecture is $p^{1/d}=q$.  Its flow is explicit:
\[
\varphi_t(x,\xi)=(x+t\nabla q(\xi),\xi).
\]
It is entire in $t$.  At the time $t=iq(\xi)$ used in Conjecture 4.1,
\[
\pi_x\varphi_{iq(\xi)}(x,\xi)
=x+iq(\xi)\nabla q(\xi)
=x+i\nabla F(\xi).
\]
Lemma~\ref{lem:legendre} therefore proves conclusion (1), with the asserted
tube, for every $\eps>0$.

The function $F^*=\tfrac12(q^*)^2$ is real analytic away from the origin, and
its Hessian is positive definite there: at dual points it is the inverse of
$D^2F$.  Thus
\[
r(x+iv)=2F^*(v)-\eps^2
\]
is a real-analytic defining function near $\partial M^\Phi_\eps$.  Its complex
Levi form is, up to a positive fixed factor, the restriction of $D^2F^*(v)$.
It is positive definite at every boundary point.  This proves conclusions
(2) and (3).

\subsection{Holomorphic continuation of the kernel}

With normalized Haar measure, the kernel of $e^{-\eps P^{1/d}}$ is
\begin{equation}\label{eq:kernel}
K_\eps(z,y)=(2\pi)^{-n}\sum_{k\in\Z^n}
e^{-\eps q(k)}e^{ik\cdot(z-y)}.
\end{equation}
Put $z=x+iv$.  Duality gives
\[
k\cdot v\ge -q(k)q^*(v),
\]
and hence
\[
\left|e^{-\eps q(k)}e^{ik\cdot(z-y)}\right|
\le e^{-(\eps-q^*(v))q(k)}.
\]
On every compact subset of $D_\eps$, the right side has uniform exponential
decay because $q(k)\asymp|k|$.  The same remains true after arbitrary
derivatives, which only add polynomial factors in $k$.  Thus
\eqref{eq:kernel} converges normally with all derivatives and is holomorphic
in $z$.  This proves conclusion (4).

\subsection{The positive complex phase and its order}

Let $\Sigma_\eps=\partial D_\eps$.  After a smooth removal of the zero
frequency, Poisson summation expresses the kernel locally, modulo a smooth
kernel, as a finite relevant lattice translate of
\begin{equation}\label{eq:osc}
(2\pi)^{-n}\int_{\R^n}
e^{i(x-y)\cdot\xi-\eps q(\xi)-v\cdot\xi}a(\xi)\,d\xi,
\qquad v\in\Sigma_\eps,
\end{equation}
where $a$ is a classical symbol of order zero equal to one for large $|\xi|$.
Equivalently, its complex phase is
\begin{equation}\label{eq:phase}
\phi(x,v,y,\xi)
=(x-y)\cdot\xi+i\bigl(\eps q(\xi)+v\cdot\xi\bigr).
\end{equation}
For $q^*(v)=\eps$,
\[
\operatorname{Im}\phi=\eps q(\xi)+v\cdot\xi\ge0.
\]
Strong convexity makes the support point unique, so equality holds exactly
when
\begin{equation}\label{eq:support}
v=-\eps\nabla q(\xi).
\end{equation}
At such a point, $\partial_\xi\phi=0$ also forces $x=y$ (locally modulo
$2\pi\Z^n$).  The mixed $(x-y,\xi)$ block is the identity, while strict
convexity is the positive transverse Hessian condition.  Hence
\eqref{eq:phase} is a nondegenerate positive complex phase in the sense of
Melin--Sj\"ostrand \cite{MelinSjostrand}.  Its canonical relation is the
graph
\begin{equation}\label{eq:graph}
(y,\xi)\longmapsto
\bigl(x=y,\ v=-\eps\nabla q(\xi);\ \xi\cdot dx\bigr).
\end{equation}
Indeed, $\xi$ is normal to $\Sigma_\eps$ at the support point, so the
$i\xi\cdot dv$ part restricts to zero on $T_v\Sigma_\eps$.

There are $N=n$ phase variables, the target has dimension $2n-1$, the source
has dimension $n$, and the amplitude has order zero.  The distributional FIO
order is therefore
\[
0+\frac{N}{2}-\frac{(2n-1)+n}{4}
=-\frac{n-1}{4}.
\]
The amplitude and phase are classical homogeneous away from the zero section,
so the operator is polyhomogeneous.  This is conclusion (5), using the
standard positive complex-phase FIO calculus
\cite{HormanderFIO,MelinSjostrand}.

\subsection{Ellipticity and the exact Sobolev shift}

The order and ellipticity can also be seen directly from the Gram operator.
Translation invariance diagonalizes $S_\eps^*S_\eps$ in the $e_k$ basis.  For
any smooth positive density $b(v)d\sigma(v)$ on $\Sigma_\eps$, its eigenvalue
is
\begin{equation}\label{eq:gram}
A_\eps(k)=e^{-2\eps q(k)}
\int_{\Sigma_\eps}e^{-2k\cdot v}b(v)\,d\sigma(v)>0.
\end{equation}
For $k=\lambda\theta$, $\lambda\to\infty$, the exponent has a unique maximum
at $v=-\eps\nabla q(\theta)$.  Its tangential Hessian is negative definite by
strong convexity.  Uniform Laplace asymptotics on the compact direction sphere
therefore give the full classical expansion
\begin{equation}\label{eq:laplace}
A_\eps(k)\sim |k|^{-(n-1)/2}
\left(a_0\!\left(\frac{k}{|k|}\right)
+\frac{a_1(k/|k|)}{|k|}+\cdots\right),
\qquad a_0>0.
\end{equation}
Thus $S_\eps^*S_\eps$ is elliptic of order $-(n-1)/2$.

Choose a positive, translation-invariant elliptic operator $\Lambda$ of order
one on $\T^n\times\Sigma_\eps$ to define the boundary Sobolev norms.  The graph
FIO composition calculus applied to
\[
B_r=S_\eps^*\Lambda^{2r}S_\eps
\]
shows that $B_r$ is a positive elliptic pseudodifferential operator of order
\[
2r-\frac{n-1}{2}.
\]
The principal symbol is positive because the graph \eqref{eq:graph} is
injective and $a_0>0$.  With
$r=s+(n-1)/4$, the order is $2s$.  Since all Fourier modes have strictly
positive norm under $S_\eps$ and the high-frequency symbol is uniformly
elliptic, the finitely many low modes cause no loss.  Consequently
\begin{equation}\label{eq:normeq}
c_s\|f\|_{H^s(\T^n)}
\le \|S_\eps f\|_{H^{s+(n-1)/4}(\partial M^\Phi_\eps)}
\le C_s\|f\|_{H^s(\T^n)}.
\end{equation}
In particular the range is closed.

It remains only to identify the range.  Every finite Fourier sum
\[
u(z)=\sum_{k\in E}c_ke^{ik\cdot z}
\]
is entire on $\T^n_\C$ and lies in the range, since the nonzero factor
$e^{-\eps q(k)}$ can be absorbed into its input coefficient.  Conversely, if
$u$ is holomorphic in $M^\Phi_\eps$ and smooth on its closure, its Fourier
coefficients in $x$ satisfy the Cauchy--Riemann equations
\[
u_k(v)=c_ke^{-k\cdot v}.
\]
Smoothness makes its Fourier truncations converge in every boundary Sobolev
norm.  Hence Fourier polynomials are dense in
$\cO^r(\partial M^\Phi_\eps)$ by the definition of that space.  The range of
$S_\eps$ is both closed and contains this dense set, so it equals
$\cO^r$.  Together with \eqref{eq:normeq}, this proves conclusion (6) and
Theorem~\ref{thm:main}.

\subsection{The explicit quartic family}

For
\[
p_\delta(\xi)=|\xi|^4+\delta\xi_1^2\xi_2^2,
\qquad q_\delta=p_\delta^{1/4},
\]
the function
$F_\delta=\tfrac12p_\delta^{1/2}$ is two-homogeneous and real analytic away
from zero.  At $\delta=0$, $D^2F_0=I$.  Hessian positivity therefore persists
for all sufficiently small $|\delta|$ by continuity on the compact unit
circle and zero-homogeneity of the Hessian.  Taking $0<\delta<\delta_0$ also
makes $p_\delta$ positive.

If $p_\delta$ were the square of a positive quadratic form
$a\xi_1^2+2b\xi_1\xi_2+c\xi_2^2$, the missing cubic monomials in its square
would force $b=0$.  The pure quartic coefficients then force $a=c=1$, and the
mixed coefficient would be $2$, contradicting the actual value $2+\delta$.
This proves Corollary~\ref{cor:quartic}.

\section{Verification notes}

\begin{enumerate}
\item \textbf{Hypotheses.}  The symbol $q^d$ is positive and homogeneous of
order $d$ away from zero.  The corresponding torus multiplier is formally
self-adjoint, elliptic, and nonnegative.  Its unit co-ball is $\{q\le1\}$,
which is strictly convex under (Q2).  The value at $k=0$ is a finite-rank,
hence smoothing, issue.
\item \textbf{Signs.}  On $z=x+iv$, $e^{ikz}=e^{ikx-kv}$.  Therefore the
boundary equality locus is the \emph{minimizing} support point
$v=-\eps\nabla q(k)$, consistent in \eqref{eq:phase}, \eqref{eq:support}, and
\eqref{eq:gram}.
\item \textbf{Order check.}  The complex phase uses $n$ variables on a kernel
whose base dimension is $(2n-1)+n$.  This independently gives
$n/2-(3n-1)/4=-(n-1)/4$.  Squaring produces order $-(n-1)/2$, matching the
$|k|^{-(n-1)/2}$ Laplace asymptotic.
\item \textbf{Sobolev check.}  Inserting a boundary operator of order $2r$
gives Gram order $2r-(n-1)/2$.  Setting $r=s+(n-1)/4$ yields exactly $2s$,
so no dimension shift is missing.
\item \textbf{No computational dependency.}  No numerical experiment or
computer algebra statement is used in the proof.
\end{enumerate}

\paragraph{Verifier report.}
The argument was checked internally by two independent order calculations:
the phase-dimension formula and the Gram/Laplace calculation.  The flow formula,
duality sign, boundary support point, and Sobolev shift agree across the two
routes.  No contradiction was found.  The most important expert-review point
is the invocation of the standard positive complex-phase FIO composition
calculus between \eqref{eq:phase} and \eqref{eq:normeq}; the explicit Gram
asymptotic supplies an independent ellipticity check but does not replace every
microlocal bookkeeping detail for arbitrary real $s$.

\section{Scope and bounded novelty check}

This is a full proof of the conjecture only for the stated flat-torus Fourier
multiplier class, so it is packaged as a partial result relative to the general
problem.  It does not treat $x$-dependent symbols, curved manifolds, lower-order
terms that do not commute with the multiplier, or loss of strong convexity.

On 21 July 2026, bounded searches were made in the run's registry, solution,
attempt, and proof-gap indexes and on arXiv/web search using the phrases
``Boutet de Monvel conjecture flat torus pseudodifferential strongly convex
Fourier multiplier'', ``Boutet de Monvel Conjecture torus'', ``holomorphic
Sobolev spaces anisotropic torus Poisson kernel'', and ``imaginary time
Hamiltonian flow Minkowski norm torus Poisson kernel''.  No exact result or
close arXiv treatment of this anisotropic flat-torus class was found.  The
source paper records only the Laplacian case as proved.  This is evidence for
novelty, not an exhaustive literature certification.

\paragraph{Human-review recommendation.}
Review as a candidate new partial theorem.  Priority checks are: (i) that the
analytic pseudodifferential convention permits the harmless smoothing choice
near the zero section; (ii) the positive complex-phase nondegeneracy and PHG
order in \eqref{eq:osc}--\eqref{eq:graph}; and (iii) the graph-composition
argument for all real Sobolev exponents.  The Legendre geometry, holomorphic
Fourier convergence, Laplace ellipticity, and explicit quartic example are
elementary and separately checkable.

\begin{thebibliography}{9}

\bibitem{Winterrose}
D. S. Winterrose,
\emph{Algebras of pseudo-differential operators acting on holomorphic Sobolev
spaces}, arXiv:2110.09389 (2021), especially Conjecture 4.1.

\bibitem{Boutet}
L. Boutet de Monvel,
\emph{Convergence dans le domaine complexe des s\'eries de fonctions propres},
Journ\'ees \'Equations aux d\'eriv\'ees partielles (1978).

\bibitem{HormanderFIO}
L. H\"ormander,
\emph{Fourier integral operators. I}, Acta Math. \textbf{127} (1971), 79--183.

\bibitem{MelinSjostrand}
A. Melin and J. Sj\"ostrand,
\emph{Fourier integral operators with complex-valued phase functions},
Lecture Notes in Math. \textbf{459}, Springer (1975).

\end{thebibliography}

\end{document}
